\documentclass{beamer}
\input{tbook-deck}
\title{Sharp bounds for weak-measurement pointer products}
\author{Haobo Ma\inst{1} \and Wenlin Zhang\inst{1,2}}
\institute{\inst{1}ChronoAI Pte Ltd, Singapore \and \inst{2}National University of Singapore}
\date{Version 2}
\begin{document}
% Manuscript: §1, §§2–3
\begin{frame}
\titlepage
\note{Ask how low an average of three signed pointer products can go. The argument moves from a rational witness to a sharp construction at every step.}
\end{frame}

% Manuscript: §§1–4
\begin{frame}{Outline}
\tableofcontents
\note{Define the observable coefficient before turning to the sharp envelope, finite strengths and evidence.}
\end{frame}

\section{Question and witness}
\subsection{Readings and order}
% Manuscript: §1, equations (1)–(2)
\begin{frame}{What does a run record?}
Each coupling has its own initially centered pointer. The two projection components shift it by $0$ or $h_j$. Keep all $n$ readings from each run.
\[
 m_n=\frac{\mathbb E[x_1\cdots x_n]}{h_1\cdots h_n},
 \qquad \mu_n=\lim_{\varepsilon\downarrow0}m_n,\quad h_j=\varepsilon g_j.
\]
Abbott et al.\ proved $\mu_2\ge-1/8$ for two projections and conjectured the same bound at all lengths.
\note{Use independent centered real Gaussian pointers. Hold $\rho$, the projections and the positive $g_j,\sigma_j$ fixed as the common scale $\varepsilon$ tends to zero. Every run enters the average.}
\end{frame}

% Manuscript: §1, equation (1)
\begin{frame}{Both orders enter the average}
Let $P(x,y)=(x,0)$ and $Q(x,y)=((x+y)/2,(x+y)/2)$. On $(0,1)$ the product $PQ$ gives $(1/2,0)$ and $QP$ gives zero.
\[
 \mathcal J_P(X)=\tfrac12(PX+XP),\quad
 X_n=P_n,\quad X_j=\mathcal J_{P_j}(X_{j+1}).
\]
\[
 \mu_n=\operatorname{Tr}(\rho X_1),\qquad \rho\succeq0,\quad \operatorname{Tr}\rho=1.
\]
\note{An orthogonal projection keeps the component in its subspace, so $P=P^*=P^2$. Hermiticity gives a real expectation. A density matrix gives a weighted average of unit pure-state expectations.}
\end{frame}

% Manuscript: §3, rational witness
\begin{frame}{A dot product in three dimensions}
$P_1$ retains the line through $(1,1,-1)$. $P_2$ and $P_3$ retain $y=z$ and $x=z$. Put $B=P_2P_3+P_3P_2$.
\[
 P_1e_3=\tfrac13(-1,-1,1)^{\mathsf T},\quad
 Be_3=(3/4,3/4,1/2)^{\mathsf T}.
\]
\begin{bookresult}{pointer-refutation}
State $e_3$ gives $\mu_3=\tfrac12\langle P_1e_3,Be_3\rangle=-1/6<-1/8$, refuting the universal bound.
\end{bookresult}
\note{The plane projections average two coordinates. The orders give $P_2P_3e_3=(1/2,1/4,1/4)^{\mathsf T}$ and $P_3P_2e_3=(1/4,1/2,1/4)^{\mathsf T}$. Symmetry turns the outer two terms into the same dot product. The Lean conclusion is the negation of the universal bound, with the equality calculated inside its proof.}
\end{frame}

\section{Sharp bounds}
\subsection{One step}
% Manuscript: §2, Theorem 1
\begin{frame}{The exact interval at each length}
\begin{bookresult}{sharp-bounds}
For finite-dimensional orthogonal projections and any density matrix,
\[
 -c_n\le\mu_n\le1,\qquad c_1=0,\quad c_{r+1}=\frac{(1+c_r)^2}{8(1-c_r)}.
\]
Real codimension-one projections and a pure state attain the lower endpoint in dimension $n$. Identities attain the upper endpoint.
\end{bookresult}
\[
 c_2=\tfrac18,\quad c_3=\tfrac{81}{448},\quad c_4=\tfrac{279841}{1315328},
 \qquad \tfrac13-c_n\sim\tfrac{16}{27n}.
\]
\note{The inequality holds for every $d,n\ge1$. The attaining construction can choose $d=n$, using the zero projection at $n=1$. Strict increase to $1/3$ gives the infimum $-1/3$ and its nonattainment at finite length.}
\end{frame}

% Manuscript: §2, equations (3)–(4)
\begin{frame}{Complete the square for the next step}
Write $X=aI+bH$ with $a=(1-c)/2$, $b=(1+c)/2$ and $\|H\|\le1$. For a unit state put $s=\|P\psi\|\in[0,1]$.
\[
 as^2-bs\le\langle\psi,\mathcal J_P(X)\psi\rangle\le as^2+bs\le1.
\]
When $0\le c\le1/3$, the vertex $s=b/(2a)$ lies in the interval. Its value is $-b^2/(4a)=-F(c)$.
\note{The midpoint and half-width of $[-c,1]$ are $a$ and $b$. Bound the remaining term $b\operatorname{Re}\langle P\psi,H\psi\rangle$ by Cauchy--Schwarz. The upper estimate uses $a+b=1$.}
\end{frame}

\subsection{Attainment and limit}
% Manuscript: §2, equations (5)–(6)
\begin{frame}{Use a fresh direction at each step}
In the plane with $X=\operatorname{diag}(1,-c)$ choose unit vectors $\psi=(\alpha,\beta)^{\mathsf T}$ and $p=(-\alpha,\beta)^{\mathsf T}$ so $\beta^2-\alpha^2=b/(2a)$. Projection onto $p$ reaches the vertex and makes Cauchy--Schwarz tight.
\[
 (e_r,v_r)\text{ plane}\quad\longrightarrow\quad
 q_r=\beta_r e_r+\alpha_r v_r,\quad P_{n-r}=I-q_rq_r^{\mathsf T}.
\]
The new bottom eigenvector is $v_{r+1}=\alpha_r e_r+\beta_r v_r$.
\note{Normalization gives $\alpha^2=(1-y)/2,\ \beta^2=(1+y)/2$ for $y=b/(2a)$. The equalities fit together because $p=-H\psi$ and $P\psi=yp$. Start from $P_n=I-e_0e_0^{\mathsf T}$. Later steps fix the unused basis directions, so a fresh $e_r$ still has eigenvalue $1$.}
\end{frame}

% Manuscript: §2, Theorem 1 and ordered product comparison
\begin{frame}{The finite-length floor tends to $-1/3$}
\[
 F(c)-c=\frac{(1-3c)^2}{8(1-c)},\qquad
 \tfrac13-F(c)=\frac{(1-3c)(c+5)}{24(1-c)}.
\]
Starting from $c_1=0$, the sequence increases strictly to the fixed point $1/3$. Abbott et al.'s ordered weak values can approach $-1$. The distance from the pointer minimum to that reference point is $1-c_n>2/3$.
\medskip

The recurrence and rate are due to Sun, Luo and Ye. The one-step Jordan estimate goes back to Strang.
\note{Distinguish the length limit from the weak-coupling limit taken at fixed length. The gap comes from the invariant interval $[-1/3,1]$ for the Jordan recursion. The appendix explains the rate.}
\end{frame}

\section{Extensions and evidence}
\subsection{Other inputs}
% Manuscript: §2, Corollaries 2 and 4
\begin{frame}{Which spectral conditions preserve bounds?}
\begin{bookresult}{positive-contractions}
For $0\preceq A_j\preceq I$, the extrema remain $-c_n$ and $1$. Expand $A_j$ as a convex combination of spectral-threshold projections.
\end{bookresult}
\begin{bookresult}{commuting-symmetric}
Commuting projections place both expectations in $[0,1]$. For Hermitian inputs $|\mu_n|\le\prod_j\|A_j\|$. With symmetric spectra, the extreme spectral products are the two endpoints of this norm interval.
\end{bookresult}
\note{The threshold weights are consecutive eigenvalue differences. The zero projection completes the weights. Multilinearity preserves the interval, with projections attaining both ends. For commuting projections both recursions yield $P_1\cdots P_n$. The norm bound uses the triangle inequality and submultiplicativity. A zero input forces the coefficient and all spectral products to vanish.}
\end{frame}

% Manuscript: §1 and §2, Corollary 3
\begin{frame}{The Gaussian envelope at positive strengths}
Let $t=e^{-h^2/(8\sigma^2)}$. The normalized first-moment map is
\[
 T_{P,t}(Y)=t\mathcal J_P(Y)+(1-t)PYP.
\]
\begin{bookresult}{gaussian-envelope}
Independent centered real Gaussian pointers satisfy $-c_n\le m_n\le1$ at all $h_j>0$. Since $T_{P_n,t_n}(I)=P_n$, the final width and coupling drop out.
\end{bookresult}
\note{Completing the square in the two shifted Gaussians gives the overlap factor. Compression preserves $[-cI,I]$ and the Jordan term lies in $[-F(c)I,I]$, so the mixture stays in the latter interval. Reducing the couplings can approach the weak-limit endpoint. The sharp minimum at prescribed positive strengths remains outside this calculation.}
\end{frame}

\subsection{Sources and scope}
% Manuscript: §4; references
\begin{frame}{Sources and scope}
Lean proves the negation of the universal $-1/8$ bound. The paper supplies written arguments for sharpness, the asymptotic and all three corollaries. The construction leaves minima at fixed $d<n$ and for all rank-one projections with $n\ge3$ outside its scope.
\smallskip
\begin{itemize}
\item Ma and Zhang, \href{https://doi.org/10.5281/zenodo.24212412}{source paper, version 2}.
\item Abbott et al., \href{https://doi.org/10.22331/q-2019-10-14-194}{Quantum 3, 194 (2019)}, §4 and Appendix C.
\item \href{https://arxiv.org/abs/1503.06387v3}{Sun, Luo and Ye, §6.2 and §6.4}. \href{https://doi.org/10.2307/2312733}{Strang, 1962}.
\end{itemize}
\note{The original question occurs immediately after equation (26) in Abbott et al. Section 3 of the paper supplies the three-dimensional witness. The result cards link the shared requests for the remaining exact statements.}
\end{frame}

\appendix
% Manuscript: §2, final paragraph of the proof of Theorem 1
\begin{frame}{Where does the $1/n$ rate come from?}
Set $\delta_n=\tfrac13-c_n$. The recurrence gives
\[
 \frac1{\delta_{n+1}}-\frac1{\delta_n}
 =\frac{27}{16-3\delta_n}\longrightarrow\frac{27}{16}.
\]
Sum the consecutive increments and divide by $n$ to get $1/(n\delta_n)\to27/16$, or $\delta_n\sim16/(27n)$.
\note{Every $\delta_n$ is positive, so the reciprocals are defined. The average of the converging increments has the same limit. The first term $1/(n\delta_1)$ vanishes.}
\end{frame}

% Manuscript: §4, coverage and verification
\FormalVerification
\end{document}
